\documentclass{article}
\usepackage{amsmath,amssymb}
\usepackage{xurl}
\usepackage[hidelinks]{hyperref}
\setlength{\emergencystretch}{2em}
% Shared commands for the notes in docs/paper-gaps/.
% Every note loads this file via % Shared commands for the notes in docs/paper-gaps/.
% Every note loads this file via % Shared commands for the notes in docs/paper-gaps/.
% Every note loads this file via \input{command}. Project-wide notation
% belongs here, never in per-note redefinitions.

% --- Project configuration ------------------------------------------------
\newcommand{\projectIssueBase}{https://github.com/LionSR/QICLean/issues/}
\newcommand{\projectPrBase}{https://github.com/LionSR/QICLean/pull/}
\newcommand{\projectDocBase}{https://sirui-lu.com/QICLean/docs/}
% ---------------------------------------------------------------------------

\DeclareUnicodeCharacter{2080}{\ifmmode{}_{0}\else\textsubscript{0}\fi}
\DeclareUnicodeCharacter{2081}{\ifmmode{}_{1}\else\textsubscript{1}\fi}
\DeclareUnicodeCharacter{2082}{\ifmmode{}_{2}\else\textsubscript{2}\fi}
\DeclareUnicodeCharacter{2099}{\ifmmode{}_{n}\else\textsubscript{n}\fi}

\newcommand{\Lean}{\textsc{Lean}}

% A formal declaration name, upright sans, breakable at underscores.
\ExplSyntaxOn
\NewDocumentCommand{\leanid}{m}{
  \ifmmode
    \textsf{\detokenize{#1}}
  \else
    \group_begin:
    \urlstyle{sf}
    \tl_set:Nn \l_tmpa_tl {#1}
    \tl_replace_all:Nnn \l_tmpa_tl {\_} {_}
    \exp_args:NV \path \l_tmpa_tl
    \group_end:
  \fi
}
\ExplSyntaxOff

% Traceability links; use in footnotes, not in the running prose.
\newcommand{\ghissue}[1]{\href{\projectIssueBase#1}{GitHub issue \##1}}
\newcommand{\ghpr}[1]{\href{\projectPrBase#1}{PR \##1}}
\newcommand{\doclink}[1]{\href{\projectDocBase#1.html}{\path{#1}}}

% Domain notation, using \providecommand so a note-local refinement stays
% possible.
\providecommand{\tr}{\operatorname{tr}}

% Dirac notation and the identity operator, as in blueprint/src/macros.
\usepackage{dsfont}
\providecommand{\ket}[1]{|#1\rangle}
\providecommand{\bra}[1]{\langle #1|}
\providecommand{\braket}[2]{\langle #1 | #2 \rangle}
\providecommand{\ketbra}[2]{|#1\rangle\!\langle #2|}
\providecommand{\Id}{\mathds{1}}
\providecommand{\one}{\mathds{1}}
\providecommand{\C}{\mathbb{C}}
\providecommand{\MN}[1]{M_{#1}(\C)}
\providecommand{\MD}{M_D(\C)}

% Verdict marker: \gapnote{<kind>}{<status>}. Kind is the policy
% classification (clarification, local-correction, scope-restriction,
% unfaithful, false-source, open-gap); status is open, wip (elimination
% underway), resolved, or historical. Machine-read by the site index;
% prints a verdict line.
\newcommand{\gapnote}[2]{\par\noindent\textbf{Verdict:} #1 (#2).\par}
. Project-wide notation
% belongs here, never in per-note redefinitions.

% --- Project configuration ------------------------------------------------
\newcommand{\projectIssueBase}{https://github.com/LionSR/QICLean/issues/}
\newcommand{\projectPrBase}{https://github.com/LionSR/QICLean/pull/}
\newcommand{\projectDocBase}{https://sirui-lu.com/QICLean/docs/}
% ---------------------------------------------------------------------------

\DeclareUnicodeCharacter{2080}{\ifmmode{}_{0}\else\textsubscript{0}\fi}
\DeclareUnicodeCharacter{2081}{\ifmmode{}_{1}\else\textsubscript{1}\fi}
\DeclareUnicodeCharacter{2082}{\ifmmode{}_{2}\else\textsubscript{2}\fi}
\DeclareUnicodeCharacter{2099}{\ifmmode{}_{n}\else\textsubscript{n}\fi}

\newcommand{\Lean}{\textsc{Lean}}

% A formal declaration name, upright sans, breakable at underscores.
\ExplSyntaxOn
\NewDocumentCommand{\leanid}{m}{
  \ifmmode
    \textsf{\detokenize{#1}}
  \else
    \group_begin:
    \urlstyle{sf}
    \tl_set:Nn \l_tmpa_tl {#1}
    \tl_replace_all:Nnn \l_tmpa_tl {\_} {_}
    \exp_args:NV \path \l_tmpa_tl
    \group_end:
  \fi
}
\ExplSyntaxOff

% Traceability links; use in footnotes, not in the running prose.
\newcommand{\ghissue}[1]{\href{\projectIssueBase#1}{GitHub issue \##1}}
\newcommand{\ghpr}[1]{\href{\projectPrBase#1}{PR \##1}}
\newcommand{\doclink}[1]{\href{\projectDocBase#1.html}{\path{#1}}}

% Domain notation, using \providecommand so a note-local refinement stays
% possible.
\providecommand{\tr}{\operatorname{tr}}

% Dirac notation and the identity operator, as in blueprint/src/macros.
\usepackage{dsfont}
\providecommand{\ket}[1]{|#1\rangle}
\providecommand{\bra}[1]{\langle #1|}
\providecommand{\braket}[2]{\langle #1 | #2 \rangle}
\providecommand{\ketbra}[2]{|#1\rangle\!\langle #2|}
\providecommand{\Id}{\mathds{1}}
\providecommand{\one}{\mathds{1}}
\providecommand{\C}{\mathbb{C}}
\providecommand{\MN}[1]{M_{#1}(\C)}
\providecommand{\MD}{M_D(\C)}

% Verdict marker: \gapnote{<kind>}{<status>}. Kind is the policy
% classification (clarification, local-correction, scope-restriction,
% unfaithful, false-source, open-gap); status is open, wip (elimination
% underway), resolved, or historical. Machine-read by the site index;
% prints a verdict line.
\newcommand{\gapnote}[2]{\par\noindent\textbf{Verdict:} #1 (#2).\par}
. Project-wide notation
% belongs here, never in per-note redefinitions.

% --- Project configuration ------------------------------------------------
\newcommand{\projectIssueBase}{https://github.com/LionSR/QICLean/issues/}
\newcommand{\projectPrBase}{https://github.com/LionSR/QICLean/pull/}
\newcommand{\projectDocBase}{https://sirui-lu.com/QICLean/docs/}
% ---------------------------------------------------------------------------

\DeclareUnicodeCharacter{2080}{\ifmmode{}_{0}\else\textsubscript{0}\fi}
\DeclareUnicodeCharacter{2081}{\ifmmode{}_{1}\else\textsubscript{1}\fi}
\DeclareUnicodeCharacter{2082}{\ifmmode{}_{2}\else\textsubscript{2}\fi}
\DeclareUnicodeCharacter{2099}{\ifmmode{}_{n}\else\textsubscript{n}\fi}

\newcommand{\Lean}{\textsc{Lean}}

% A formal declaration name, upright sans, breakable at underscores.
\ExplSyntaxOn
\NewDocumentCommand{\leanid}{m}{
  \ifmmode
    \textsf{\detokenize{#1}}
  \else
    \group_begin:
    \urlstyle{sf}
    \tl_set:Nn \l_tmpa_tl {#1}
    \tl_replace_all:Nnn \l_tmpa_tl {\_} {_}
    \exp_args:NV \path \l_tmpa_tl
    \group_end:
  \fi
}
\ExplSyntaxOff

% Traceability links; use in footnotes, not in the running prose.
\newcommand{\ghissue}[1]{\href{\projectIssueBase#1}{GitHub issue \##1}}
\newcommand{\ghpr}[1]{\href{\projectPrBase#1}{PR \##1}}
\newcommand{\doclink}[1]{\href{\projectDocBase#1.html}{\path{#1}}}

% Domain notation, using \providecommand so a note-local refinement stays
% possible.
\providecommand{\tr}{\operatorname{tr}}

% Dirac notation and the identity operator, as in blueprint/src/macros.
\usepackage{dsfont}
\providecommand{\ket}[1]{|#1\rangle}
\providecommand{\bra}[1]{\langle #1|}
\providecommand{\braket}[2]{\langle #1 | #2 \rangle}
\providecommand{\ketbra}[2]{|#1\rangle\!\langle #2|}
\providecommand{\Id}{\mathds{1}}
\providecommand{\one}{\mathds{1}}
\providecommand{\C}{\mathbb{C}}
\providecommand{\MN}[1]{M_{#1}(\C)}
\providecommand{\MD}{M_D(\C)}

% Verdict marker: \gapnote{<kind>}{<status>}. Kind is the policy
% classification (clarification, local-correction, scope-restriction,
% unfaithful, false-source, open-gap); status is open, wip (elimination
% underway), resolved, or historical. Machine-read by the site index;
% prints a verdict line.
\newcommand{\gapnote}[2]{\par\noindent\textbf{Verdict:} #1 (#2).\par}

\title{Quantum Neyman--Pearson: Freedom on the Kernel of the Optimal Test}
\author{}
\date{2026-11-15}
\begin{document}
\maketitle
\gapnote{false-source}{resolved}

\section{Source statement}
Let $\rho,\sigma$ be density matrices, let $p$ be the prior probability of
$\rho$, and let $0\le P\le\mathbb 1$ be the effect on which $\rho$ is guessed.
Wolf's Theorem ``Quantum Neyman--Pearson'' \cite{Wolf2012Quantum}%
\footnote{\path{Notes/WolfNoteTexSource/ch08_distance_measures.tex},
lines~543--566, Equations~(8.40)--(8.42).}
bounds the average error
\begin{align}
    e(P) &= p\tr[\rho(\mathbb 1-P)]+(1-p)\tr[\sigma P]
    \label{eq:np_kernel_error}
\end{align}
by
\begin{align}
    e(P) &\ge \tfrac12\bigl(1-\lVert A\rVert_1\bigr),
    \qquad A=p\rho-(1-p)\sigma,
    \label{eq:np_kernel_bound}
\end{align}
and asserts that equality holds if and only if $P$ is the projection onto
the positive part of $A$.

\section{The equality clause}
Write $A=A_+-A_-$ with $A_\pm\ge0$ and $A_+A_-=0$. For every matrix $P$,
using $\tr\rho=\tr\sigma=1$,
\begin{align}
    e(P)-\tfrac12\bigl(1-\lVert A\rVert_1\bigr)
        &= \tr[(\mathbb 1-P)A_+]+\tr[PA_-].
    \label{eq:np_kernel_excess}
\end{align}
For an effect both traces are nonnegative, which gives
(\ref{eq:np_kernel_bound}). A trace of a product of two positive
semidefinite matrices vanishes only if the product vanishes. Hence equality
holds if and only if
\begin{align}
    PA_+ &= A_+, & PA_- &= 0.
    \label{eq:np_kernel_condition}
\end{align}
This is the condition stated in Wolf's proof. It fixes $P$ on the range of
$A_+$ and on the range of $A_-$, but imposes nothing on the kernel of $A$.

\paragraph{Counterexample to the printed clause.}
Take $p=\tfrac12$ and $\rho=\sigma$. Then $A=0$, the projection onto the
positive part of $A$ is $0$, and every effect satisfies
\begin{align}
    e(P) &= \tfrac12\tr[\rho(\mathbb 1-P)]+\tfrac12\tr[\rho P]
          = \tfrac12
          = \tfrac12\bigl(1-\lVert A\rVert_1\bigr).
    \label{eq:np_kernel_counterexample}
\end{align}
So $P=\mathbb 1$ is optimal but is not the projection onto the positive part.

\paragraph{Invertible difference.}
If $A$ is invertible, (\ref{eq:np_kernel_condition}) gives
$PA=A_+=\Pi_+A$, where $\Pi_+$ is the projection onto the range of $A_+$.
Cancelling $A$ yields $P=\Pi_+$. Thus the printed clause holds exactly when
the kernel of $A$ is trivial.

\section{Formalization}
\doclink{QICLean/Analysis/Optimization/QuantumNeymanPearson} proves
(\ref{eq:np_kernel_excess}), the bound (\ref{eq:np_kernel_bound}), the
corrected equivalence with (\ref{eq:np_kernel_condition}), attainment by
$\Pi_+$, the printed clause for invertible $A$, and the counterexample
(\ref{eq:np_kernel_counterexample}). The states need only be Hermitian of
unit trace and $p$ may be any real number, which contains the source
hypotheses. The theorem for invertible $A$ carries a scope marker; the other
results are faithful to the source with the corrected equality clause. This
resolves the Neyman--Pearson part of \ghissue{46}.

\section{Verdict}
The equality clause is false as printed whenever the kernel of
$p\rho-(1-p)\sigma$ is nonzero. The correct clause is
(\ref{eq:np_kernel_condition}), which the source proof already derives.

\bibliographystyle{plain}
\bibliography{references}
\end{document}
